\documentclass{beamer}
\usepackage{subfig}
\usepackage{multicol}

\usetheme[subsectionpage=progressbar]{metropolis}

\titlegraphic{%
    \hfill
    \includegraphics[width=.4\textwidth]{../logo-uniurb-2016.jpg}
}

\setbeamertemplate{section in toc}[sections numbered]
\setbeamertemplate{subsection in toc}[subsections numbered]

\makeatletter
\setbeamertemplate{title page}{
  \begin{minipage}[b][\paperheight]{\textwidth}
    \vfill%
    \ifx\inserttitle\@empty\else\usebeamertemplate*{title}\fi
    \ifx\insertsubtitle\@empty\else\usebeamertemplate*{subtitle}\fi
    \usebeamertemplate*{title separator}
    \ifx\beamer@shortauthor\@empty\else\usebeamertemplate*{author}\fi
    \ifx\insertinstitute\@empty\else\usebeamertemplate*{institute}\fi
    \vspace*{.3cm}
    \ifx\insertdate\@empty\else\usebeamertemplate*{date}\fi
    \vfill
    \ifx\inserttitlegraphic\@empty\else\inserttitlegraphic\fi
    \vspace*{1cm}
  \end{minipage}
}

\makeatother
%-----------------------------------------------
%--------------- Copertina ---------------------
%-----------------------------------------------
\title{Progettazione e sviluppo di un'applicazione per analisi stabilometrica tramite smartphone}
\author{Lorenzo Calisti}
\institute{Corso di Laurea in Informatica Applicata \\ Anno accademico 2019-2020}
\date{26-06-2020}

%-----------------------------------------------
%---------------- Slides -----------------------
%-----------------------------------------------
\begin{document}
    \begin{frame}
        \maketitle
    \end{frame}
    
    \begin{frame}{Indice}
        \begin{multicols}{2}
        %\tableofcontents[hideallsubsections]   
        \tableofcontents
        \end{multicols}
    \end{frame}

    \begin{frame}{Download}
        \centering
        \includegraphics[scale=.6]{../figures/qrcode.png}
        \vspace{1em}
        \\
        {\itshape https://play.google.com/store/apps/details?id=srl.digit.balance}
    \end{frame}

    \section{L'analisi dell'equilibrio}
    %\section{Le pedane di forza}
    \begin{frame}{Le pedane di forza}
        \centering
        \includegraphics[scale=0.45]{../figures/pedana.png}
    \end{frame}
    \begin{frame}{Le pedane di forza}
        \centering
        \includegraphics[scale=0.3]{../figures/statokinesigramma.jpg}
    \end{frame}

    \subsection{Il pendolo inverso}
    \begin{frame}{Il pendolo inverso}
        \centering
        \includegraphics[scale=.35]{../figures/single_inverted_pendulum.png}
    \end{frame}

    \subsection{Le features}
    \begin{frame}{Le features}
        \begin{itemize}
            \item Features nel dominio del tempo
            \pause
            \item Features nel dominio della frequenza
            \pause
            \item Features strutturali
            \pause
            \item Features giroscopiche
        \end{itemize}
    \end{frame}

    \section{I sensori negli smartphone}
    \begin{frame}{I sensori negli smartphone}
        \begin{itemize}
            \item accelerometro
            \pause
            \item giroscopio
            \pause
            \item gps
            \pause
            \item sensori di prossimità
            \pause
            \item lettore di impronte digitali
            \pause
            \item sensore d'umidità
            \pause
            \item termometro
            \pause
            \item cardiofrequenzimetro
        \end{itemize}
    \end{frame}
    
    \begin{frame}{Accelerometro e Giroscopio}
        \begin{figure}
            \centering
            \subfloat{\includegraphics[scale=0.266]{../figures/accelerometer.jpg}}
            \subfloat{\includegraphics[scale=0.2]{../figures/gyroscope.jpg}}
        \end{figure}
    \end{frame}

    \section{Precedente sviluppo in Android}
    \begin{frame}{Precedente sviluppo in Android}
        \centering
        \includegraphics[scale=0.08]{../figures/android_kotlin.png}
    \end{frame}

    \section{Mockup con Figma}
    \begin{frame}
        \begin{figure}
            \centering
            \includegraphics[width=\textwidth]{../figures/figma.png}
        \end{figure}
    \end{frame}
    
    \section{L'applicazione}
    \begin{frame}
        \centering
        \includegraphics[scale=0.17]{../figures/flutter_dart.png}
    \end{frame}

    \subsection{Il Test della postura}
    \begin{frame}{Il test della postura}
        \centering
        \includegraphics[scale=0.09]{../figures/screenshot/redmi_note_8t/home.png}
    \end{frame}
    \begin{frame}{Il test della postura}
        \centering
        \includegraphics[scale=0.09]{../figures/screenshot/redmi_note_8t/home_measure.png}
    \end{frame}
    \begin{frame}{Il test della postura}
        \centering
        \includegraphics[scale=0.09]{../figures/screenshot/redmi_note_8t/home_measuring.png}
    \end{frame}

    \subsection{Plugin per la lettura dei sensori}
    \begin{frame}{Plugin per la lettura dei sensori}
        Una libreria di codice in grado di chiamare API specifiche della piattaforma per {\bfseries ottenere più prestazioni} e {\bfseries utilizzare funzionalità specifiche della piattaforma}.
    \end{frame}
    \begin{frame}{Plugin per la lettura dei sensori}
        \begin{enumerate}
            \item configurare il tempo di lettura a priori
            \pause
            \item annullare la lettura prima del termine
            \pause
            \item ottenere una sequenza con frequenza di 100Hz
        \end{enumerate}
    \end{frame}

    \subsection{Il Flutter Isolate}
    \begin{frame}{Il Flutter Isolate}
        \centering
        \includegraphics[scale=0.5]{../figures/isolate.png}
    \end{frame}

    \subsection{Dati ottenuti dall'analisi}
    \begin{frame}{Dati ottenuti dall'analisi}
        \begin{figure}
            \subfloat{\includegraphics[scale=0.09]{../figures/screenshot/redmi_note_8t/result.png}}
            \hspace{3em}
            \subfloat{\includegraphics[scale=0.09]{../figures/screenshot/redmi_note_8t/result_bottom.png}}
        \end{figure}
    \end{frame}

    \subsection{Dati d'anamnesi}
    \begin{frame}{Dati d'anamnesi}
        \begin{itemize}
            \item altezza
            \pause
            \item età
            \pause
            \item genere
            \pause
            \item peso
            \pause
            \item problemi di postura
            \pause
            \item problemi di postura in famiglia
            \pause
            \item uso di farmaci che interferiscono con la postura
            \pause
            \item altri traumi
            \pause
            \item difetti visivi/uditivi
        \end{itemize}
    \end{frame}
    \begin{frame}{Dati d'anamnesi}
        \begin{figure}
            \subfloat{\includegraphics[scale=0.09]{../figures/screenshot/redmi_note_8t/height.png}}
            \hspace{0.3em}
            \subfloat{\includegraphics[scale=0.09]{../figures/screenshot/redmi_note_8t/general_info.png}}
            \hspace{0.3em}
            \subfloat{\includegraphics[scale=0.09]{../figures/screenshot/redmi_note_8t/posture.jpg}}
        \end{figure}
    \end{frame}

    \section{Sviluppi futuri}
    \begin{frame}{Sviluppi futuri}
        \begin{itemize}
            \item implementazione e il testing in IOS
            \pause
            \item implementazione di un database remoto
            \pause
            \item miglioramento dei dati d'anamnesi
        \end{itemize}
    \end{frame}

    \section{Domande?}
\end{document}